\section{The complete local height-value set away from five}
\label{sec:local-height-values}

Let \(C\) be the smooth projective curve
\[
 W^2=z^6-27z^4+99z^2-9.
\]
The standard elliptic quotients in
\cite[Theorem~1.4]{BalakrishnanDograQCI} are
\begin{align}
 F_1&:Y^2=X^3-27X^2+99X-9,&
 f_1(z,W)&=(z^2,W),\\
 F_2&:Y^2=X^3+99X^2+243X+81,&
 f_2(z,W)&=(-9/z^2,-9W/z^3).
 \label{lh-standard-maps}
\end{align}
The maps extend to the smooth projective curves. Their globally minimal
models and corresponding points are
\begin{align}
 E&:y^2=x^3-9x-9,&
 R_1&=((z^2-9)/4,W/8),\\
 E'&:y^2=x^3-189x+999,&
 R_2&=((33-9/z^2)/4,-9W/(8z^3)).
 \label{lh-minimal-maps}
\end{align}
Both isomorphisms have scale \(u=2\), although their translations differ.
Their minimal discriminants are \(2^4 3^6\) and \(2^4 3^{10}\).

For a finite prime \(\ell\ne5\), put
\(\chi_\ell(a)=-v_\ell(a)\log_5(\ell)\), where \(\log_5(5)=0\).
Use local N\'eron heights \(h_{E,\ell},h_{E',\ell}\) in twice
Silverman's normalization on these minimal models. The local-height
coefficient is the ordinary rational coefficient of the real local
height, with \(\log(\ell)\) replaced by \(\log_5(\ell)\).
In particular, at a point of nonsingular reduction it is
\[
 h_{E,\ell}(R)=\max(0,-v_\ell(x(R)))\log_5(\ell).
\]
At singular reduction we will use precisely the valuation formula stated
below.

On the standard \(F_i\) models use the corresponding Weierstrass
tangential normalization. The exact change-of-variables rule is
\begin{equation}
 h_{F_i,\ell}(f_i(P))
   =h_{E_i,\ell}(R_i(P))+2\chi_\ell(2),
 \qquad E_1=E,\quad E_2=E'.
 \label{lh-local-transport}
\end{equation}
Indeed, the normalized local N\'eron function for \(2O\) has principal
term \(-2\chi_\ell(t)\), with \(t=-x/y\). Under
\(X=u^2x+r,\ Y=u^3y\), one has
\(t_F=t_E/u+O(t_E^3)\). Thus its tangential normalization changes by
\(2\chi_\ell(u)\), and uniqueness with this normalization proves
\eqref{lh-local-transport}. This also agrees with the local addition
identity of \cite[Lemma~7.4]{BalakrishnanDograQCI}, since
\(x_F(P)-x_F(Q)=u^2(x_E(P)-x_E(Q))\).
The global scalar is unchanged because \(\sum_v\chi_v(u)=0\);
the local functions individually have the stated constants.

Those two constants are equal and cancel. Hence the exact combination
used in Theorem~1.4 of that source is
\begin{align}
 J_\ell(P)
 &=h_{F_1,\ell}(f_1(P))-h_{F_2,\ell}(f_2(P))
                                  -2\chi_\ell(z)\notag\\
 &=h_{E,\ell}(R_1(P))-h_{E',\ell}(R_2(P))
                                  -2\chi_\ell(z).
 \label{lh-combination}
\end{align}
Initially \(z\) is finite and nonzero. The locally constant extensions
over zero and infinity are included in the proofs; we do not insert a
separate infinite height of an elliptic origin into this expression.

We recall the two cases of the local theorem that suffice here.
For a minimal Weierstrass model let
\begin{align*}
 A&=v_\ell(3x^2+2a_2x+a_4-a_1y),\\
 B&=v_\ell(2y+a_1x+a_3),\\
 C_0&=v_\ell(3x^4+b_2x^3+3b_4x^2+3b_6x+b_8).
\end{align*}
If \(A\le0\) or \(B\le0\), the height coefficient is
\(\max(0,-v_\ell(x))\). Otherwise, if \(v_\ell(c_4)>0\) and
\(C_0\ge3B\), the coefficient is \(-2B/3\).
These are cases (a) and (c) of
\cite[Proposition~3.4.1, printed pp.~72--73]{Cremona1997Algorithms}.
The text explicitly states its applicability to points over
\(\mathbb Q_\ell\) and the twice-Silverman normalization.
We use \(v_\ell(0)=+\infty\). No inference from a software Kodaira
type or a finite table is required.

\begin{theorem}[The full three-adic value set]
\label{lh-three}
For every \(P\in C(\mathbb Q_3)\), with the extended convention above,
\[
 J_3(P)=\frac43\log_5(3).
\]
Thus the full local value set is the displayed singleton.
\end{theorem}
\begin{proof}
First suppose \(z\ne0,\infty\) and let \(m=v_3(z)\).
If \(m>0\), the right side of the sextic has valuation exactly two
and nonsquare unit part \(-1\) modulo three: the other three
valuations are \(6m,3+4m,2+2m\), all larger than two.
At \(z=0\) the same obstruction is \(W^2=-9\).
Thus every finite point has \(m\le0\). If \(m=0\), then \(W\) is
a unit; if \(m<0\), the leading term is uniquely dominant, so
\(v_3(W)=3m\).

For the second minimal point \((x,y)=R_2(P)\), both cases give
\begin{equation}
 v_3(x)=1,\qquad v_3(y)=2,\qquad
 X:=x/3=\frac{11}4-\frac3{4z^2}\equiv2\pmod3.
 \label{lh-three-coordinates}
\end{equation}
For \(E'\) the third division polynomial has the identities
\begin{align}
 \psi_3(x)&=3x^4-1134x^2+11988x-35721,\\
 \psi_3(3X)&=3^5(X^4+148X)-3^6(14X^2+49).
 \label{lh-division-polynomial}
\end{align}
Since \(X\equiv2\pmod3\), one has \(X^4+148X\equiv0\pmod3\);
hence \(C_0\ge6\). Also
\[
 B=2,\qquad A=v_3(27(X^2-7))>0,\qquad v_3(c_4(E'))=4>0.
\]
The second stated case of the local theorem therefore applies and gives
\[
 h_{E',3}(R_2)=-\frac43\log_5(3).
\]
For \(E\), when \(m=0\) the ordinate \(W/8\) is a unit, so case (a)
gives zero. When \(m<0\), the abscissa has valuation \(2m\) and
\(v_3(3x^2-9)=1+4m<0\), giving
\(h_{E,3}(R_1)=-2m\log_5(3)\).
Thus that latter formula holds for both ranges \(m\le0\).
The term \(-2\chi_3(z)=2m\log_5(3)\) cancels it, proving the result
on the finite chart.

There are two rational points at infinity because the sextic is monic.
Every sufficiently close punctured three-adic neighborhood has \(m<0\),
where the combination is the same constant. It extends with that value
at both infinity points. There is no zero-fiber point over
\(\mathbb Q_3\), as already shown. This exhausts the smooth projective
local curve. Finally \((z,W)=(1,8)\) is a rational point, proving
nonemptiness and therefore equality of the value set with the singleton.
\end{proof}

\begin{theorem}[Two and all remaining primes away from five]
\label{lh-other-primes}
For every finite prime \(\ell\ne3,5\) and every
\(P\in C(\mathbb Q_\ell)\), the extended combination is
\[
 J_\ell(P)=0.
\]
\end{theorem}
\begin{proof}
At two, put \(m=v_2(z)\) on the finite nonzero chart.
If \(m<0\), the two abscissae in \eqref{lh-minimal-maps} have
valuations \(2m-2\) and \(-2\). If \(m>0\), their valuations are
\(-2\) and \(-2m-2\). For either short model and any negative
abscissa valuation, the derivative \(3x^2+a_4\) has valuation
\(2v_2(x)<0\). The local theorem gives height
\(-v_2(x)\log_5(2)\). Subtracting the two heights and adding
\(2m\log_5(2)\) gives zero in both ranges.

If \(m=0\), then \(z^2\equiv1\pmod8\). Both abscissae are integral
and even, since \(z^2-9\equiv0\pmod8\) and
\(33-9/z^2\equiv0\pmod8\). The two cubic equations then force the
ordinates to be integral odd units. The derivative \(3x^2+a_4\)
is odd, so both heights vanish; this also covers an abscissa equal
to zero. The character term vanishes as well. Thus every valuation
range gives \(J_2=0\).
The fiber \(z=0\) has no \(\mathbb Q_2\)-point because \(-1\) is
not a square in \(\mathbb Q_2\). At either infinity point the
constant punctured neighborhoods with \(m<0\) give the extension.

Now let \(\ell>3\), \(\ell\ne5\). Both minimal elliptic curves have
good reduction. With \(m=v_\ell(z)\), if \(m<0\), \(R_1\) has
abscissa valuation \(2m\) and \(R_2\) is integral. If \(m>0\),
\(R_1\) is integral and \(R_2\) has abscissa valuation \(-2m\).
For \(m=0\) both are integral. Possible divisibility of the constant
33 causes no exception: when \(-9/z^2\) is dominant its coefficient
is a unit, and otherwise extra cancellation leaves an integral
abscissa. The good-reduction height formula and the term
\(2m\log_5(\ell)\) again cancel exactly in all three cases.

At any local point above \(z=0\) or infinity, sufficiently close points
outside that fiber have respectively \(m>0\) or \(m<0\), so the
combination extends as zero. This covers every projective exception.
The rational point \((1,8)\) ensures nonemptiness at every prime.
\end{proof}

\begin{corollary}[The exact target value in quadratic Chabauty]
\label{lh-omega}
With the stated character and height conventions, the set \(\Omega\)
of \cite[Theorem~1.4]{BalakrishnanDograQCI} for this curve and \(p=5\)
is exactly
\[
 \Omega=\left\{-\frac43\log_5(3)\right\}.
\]
\end{corollary}
\begin{proof}
That theorem defines \(\Omega\) as the values of
\(-\sum_{\ell\ne5}J_\ell(P_\ell)\) on the product of the local curves.
Theorems~\ref{lh-three} and \ref{lh-other-primes} evaluate every
summand on its entire local domain, and their domains are nonempty.
Only the summand at three is nonzero, with the value asserted above.
\end{proof}

For the next calculation, the minimal-model expression on the finite
nonzero five-adic chart can be written
\begin{align}
 \mathcal F(P)
  ={}&h_{E,5}(R_1)-h_{E',5}(R_2)-2\log_5(z)\notag\\
    &-\alpha_E\log_E(R_1)^2
      +\alpha_{E'}\log_{E'}(R_2)^2.
 \label{lh-five-equation}
\end{align}
The equal local model constants cancel. Moreover the exact scaling of
Section~\ref{sec:height-normalization} gives
\(\alpha_F\log_F^2=\alpha_E\log_E^2\).
Thus this is the standard expression, with its models and signs fixed.
Every rational point in that chart satisfies
\[
 \mathcal F(P)=-\frac43\log_5(3).
\]
One can independently check the sign by summing the local heights
and using the rank-one global quadratic identities.
Since \(b=(1,8)\) is rational, the same equation is
\(\mathcal F(P)-\mathcal F(b)=0\). This relative form cancels a
uniform additive convention if the same convention is used at both
points.

No analytic zero list or multiplicity calculation is claimed here.
The remaining work is to construct compatible Coleman or sigma
expressions, prove complete truncation bounds on all required
five-adic charts, certify every analytic zero, and apply the rational
sieve. The zero fiber is over \(\mathbb Q(i)\); it has
\(\mathbb Q_5\)-points but no rational points. The two infinity points
are rational and must be recorded separately if omitted from a chart.
In particular the translated formula cannot simply assume that
\(Q_1=(0,3i)\in F_1\) is rational over \(\mathbb Q\); the direct
Theorem~1.4 expression was used above.
After a classification of \(C(\mathbb Q)\), the second square and
positive common-source conditions for the original curve \(D\)
remain additional tests. Moving exponents, residuals and the uniform
point-height problem are unchanged. These universal local-value
statements are ordinary proofs with stated standard inputs, not
a numerical sampling result or a new Lean formalization.
